\section{Cusp lifting on the ordinary locus}
\label{sec:cusp-lifting}

We pass from the theta sections of
\cref{prop:theta-family,prop:constant-term-small,prop:bounded-content}
to the Hecke algebra acting on classical cusp forms.  Throughout this
section, ordinary refers to the integral PEL variety for the target
unitary group.  Its integral datum is hyperspecial and totally split at
$p$; the argument places no reduction condition on $A_0$.

There are two steps.  We first lift the theta sections, whose constant
terms vanish modulo increasing precision, to genuine characteristic-zero
cusp forms in a higher weight.  We then use their bounded nonconstant
content to show that every relation among the cusp-form Hecke operators
vanishes at the theta comparison values.  The second step produces a
homomorphism from the entire operator order to the truncated tame ring.
Both steps use the same integral coefficient lattice.

Put $b=b_*$ and $\cO_b=\cO[t]/(t^b)$.  We use the following combined
output of the preceding sections.  After one fixed clearing of
denominators, the theta sections have bounded coefficients in the
dual-Fock lattice of \cref{eq:dual-fock-lattice}, expressed in
ordinary differential frames, and their constant terms tend to zero
coefficientwise in $\cO_b$.  At augmentation, one nonconstant
Fourier--Jacobi coefficient, evaluated in a unit frame, has valuation
bounded above by a constant independent of $i$.  Their weight at each
real place is $(m,m,1;-1)$, and $U_2$ and $U_3$ have eigenvalue $p^2$
at each distinguished $p$-adic place.  All the fixed finite local
conditions of \cref{prop:bounded-content} are retained.  In particular,
these include the binary controls at the indicated $(2,2)$ levels and
the $(3,1)$ parahoric condition at the indicated Steinberg places.
The coefficient fields of individual sections and evaluations may
vary with $i$; all valuation bounds below are uniform under these
finite coefficient extensions.  We normalize $v_p(p)=1$.

\subsection{The canonical branch of the actual level}

Choose a neat tame base level.  The depth $s$ of the $p$-level is
fixed.  Enlarge $k$ once to contain the required $p^s$-th roots of
unity and choose compatible cyclotomic frames.  All models and
finite normalizations in this section are taken after this fixed
base extension.  Write $X^{\mathrm{tor}}$ and
$X^{\min}$ for compatible hyperspecial integral toroidal and minimal
compactifications.  We use the canonical differential bundles, the
formal toroidal charts over abelian extension-data schemes, and the
ample total Hodge line bundle $\mathcal L$ on $X^{\min}$.  These are
the good-prime PEL constructions of \citet{Lan2013}; the formal-chart
statement used here is Theorem~6.4.1.1, and the Fourier--Jacobi
descriptions are in Section~7.1.2, of the author revision specified in
that reference.  The additional $p$-level will be imposed by finite
normalization, rather than by an assumption that this finer level is
hyperspecial.

For the corresponding ordinary-level constructions, see
\citet[Definition~3.2.2.9 and Propositions~5.2.3.3,
5.2.3.18]{Lan2018Ordinary}.  We identify the particular branch and
prove the base-change statement needed here below.

Let $\pi:Y^{\mathrm{tor}}\to X^{\mathrm{tor}}$ be the finite
normalization in the generic-fiber level $J$.  At each distinguished
$p$-adic factor the ordinary $p$-divisible group has multiplicative
height three and etale height one.  The level $J$ has the upper
$(2,1,1)$ block structure: the first rank-two Levi is congruent to the
identity at the chosen fixed depth, the last two unit factors and the
upper radical are full, and the lower blocks have the prescribed
depths.  Equivalently, it specifies a rank-two flag with the indicated
frame inside a rank-three flag.  It does not mark lifts of the etale
quotient.

\begin{lemma}[Canonical ordinary level]
\label{lem:canonical-ordinary-level}
Over the ordinary formal completion, the locus on which the
rank-three flag equals the multiplicative subgroup is an open and
closed branch of $Y^{\mathrm{tor}}$.  Its residual flag and frame
data are finite etale data on the Cartier-dual character group.
On a toroidal cusp chart this branch adds finite etale covers of the
abelian extension-data scheme and does not ramify the center Fourier
variable.
\end{lemma}

\begin{proof}
The ordinary group has a canonical multiplicative subgroup at every
fixed level $p^s$
\citep[proof of Theorem~2.1, pp.~149--150]{Katz1981SerreTate}.
Once the rank-three flag is this subgroup, its
rank-two subflag and frame are equivalent by Cartier duality to
quotients and markings of its finite etale character group.  The full
upper radical of $J$ is essential: specifying a splitting of the etale
quotient would impose different data.  The canonical subgroup
separates the corresponding generic level choices.  Their idempotents
extend in the finite normalization over the ordinary completion;
equivalently, one may normalize after completion, since the base is
excellent.  This proves the asserted branch decomposition.

At a rank-one boundary component, the Raynaud semiabelian group is
an extension of the binary abelian group by a rank-one torus.  The
preimage of the binary multiplicative torsion is an extension of
multiplicative-type finite flat groups, hence is of multiplicative
type: its Cartier dual is an extension of finite etale groups.  This
is the entire distinguished rank-three multiplicative subgroup.  Its
construction precedes quotienting by the Raynaud period lattice and
does not involve the center Fourier parameter.  Its character-group
level therefore produces only the asserted finite etale chart
covers.

After passage to a completed strict boundary base and choice of an
origin, a connected finite etale cover of an abelian chart is an
abelian scheme and the covering map is an isogeny.  On a geometric
fiber this follows by lifting addition to the pointed cover: the
subgroup defining the cover in the abelian fundamental group is
preserved by addition, and uniqueness of pointed lifts gives the
group laws.  The same construction extends over the strict base by
the equivalence for finite etale covers under henselian thickening.
Thus the new charts have the same abelian geometry needed for
Fourier--Jacobi coefficients; compare the ordinary chart description
in \citet[Proposition~4.2.1.34 and Corollary~4.2.1.36]{Lan2018Ordinary}.
\end{proof}

We will use ordinary expansions to test vertical integrality.
Every component at the neat hyperspecial base level has an ordinary
cusp: the boundary datum is split ordinary CM, and component
comparison for the smooth good-prime compactification identifies
these components with those in characteristic zero.  The ordinary
locus is dense.  Each canonical finite etale component also meets
the corresponding completed cusp charts, since a connected finite
etale cover maps onto its base component.  Consequently vanishing
of all reduced expansions detects a vertical factor of a section;
successive reduction gives the same assertion modulo every power
of a uniformizer.

The expansion representatives can be chosen integral at $p$, at
$q_i$, and at all the fixed exceptional places simultaneously.
Indeed, Iwasawa decomposition reduces this to weak approximation in
the cusp parabolic; the binary unitary factor has a rational Cayley
chart.  The same approximation permits the prime-to-$p$ Jacobi
torsions used for bounded content, with denominators prime to the
other specified places.  At the remaining places spherical Iwasawa
decomposition applies.  The unipotent changes do not alter a
constant term, and the Levi changes translate its binary argument
or multiply its frame by a unit at $p$.  Thus the bounds from the
preceding sections apply in all these charts.

\subsection{Positive Fourier--Jacobi coefficients and base change}

Let $D$ be the reduced toroidal boundary of $X^{\mathrm{tor}}$.
For the weights in use, put $d=m-1$ and retain the coefficient
functor from \cref{eq:dual-fock-lattice}:
\[
 \mathcal D_d(\mathcal H)
   =\bigl(\operatorname{Sym}^d
          ((\textstyle\bigwedge^2\mathcal H)^\vee)\bigr)^\vee.
\]
Here $\mathcal H=\omega^+$ is the rank-three positive Hodge bundle.
The inverse Fock-frame convention fixed there identifies this
functor with the literal dual of the Fock monomial lattice.  Define
$\mathcal E_\lambda$ on $Y^{\mathrm{tor}}$ using, at each
distinguished embedding, the positive factor
\[
       (\det\omega^+)^{\lambda_3}\otimes\mathcal D_d(\omega^+)
\]
and the opposite differential line to exponent $-\lambda_4$.
In particular this is the same integral coefficient bundle as in
the theta expansions; a Hasse shift changes only its two line-bundle
factors.  Define the coherent sheaf
\[
 \mathcal M_\lambda
 = (Y^{\mathrm{tor}}\longrightarrow X^{\min})_*
   \bigl(\mathcal E_\lambda(-\pi^*D)\bigr).
\]
The pulled-back boundary is used in this definition.  By
\cref{lem:canonical-ordinary-level}, its restriction to the canonical
ordinary cusp charts has the usual integral Fourier indices.

\begin{lemma}[Integral operations on the coefficient lattice]
\label{lem:dual-fock-functor}
The functor $\mathcal D_d$ is finite locally free and commutes with
base change.  Integral changes of Hodge frame preserve its lattice.
If a map on $\mathcal H$ is diagonal with entries $c_1,c_2,c_3$,
its action on the basis dual to the minor monomials has entries
\[
       (c_1c_2)^{a_{12}}(c_1c_3)^{a_{13}}
          (c_2c_3)^{a_{23}},
       \qquad a_{12}+a_{13}+a_{23}=d.
\]
A filtration of $\mathcal H$ by subbundles with constant locally
free graded pieces induces such a filtration on
$\mathcal D_d(\mathcal H)$.
\end{lemma}

\begin{proof}
For a finite locally free module $\mathcal W$, its divided-power
module is
\[
       \Gamma^d(\mathcal W)
          =\bigl(\operatorname{Sym}^d(\mathcal W^\vee)\bigr)^\vee.
\]
Use this integral definition with
$\mathcal W=\bigwedge^2\mathcal H$.  A minor monomial and its dual
basis vector have pairing one, without a factorial.  All the modules
in this definition are finite locally free, so it commutes with
base change.  Functoriality for linear maps proves preservation
under integral frame changes and gives the displayed diagonal
entries, including when some $c_j$ is not a unit.

To check the filtration assertion, split the filtration locally.
Exterior and symmetric powers have their degree filtrations, with
graded pieces given by tensor products of the corresponding powers
of the original graded pieces.  Dualizing the symmetric-power
filtration gives the divided-power filtration.  Block-triangular
changes between local splittings preserve these filtrations, so
they descend.  Their graded pieces are constant whenever those of
$\mathcal H$ are constant.  No rational splitting or rational
identification with ordinary symmetric powers is used.
\end{proof}

\begin{lemma}[Cusp base change on the canonical branch]
\label{lem:cusp-base-change}
An ordinary section modulo $p^M$ on the canonical branch whose
constant terms vanish, extended by zero on the other branches,
is a section of $\mathcal M_\lambda/p^M$ over the ordinary open.
The assertion remains true with coefficients in
$\cO[t]/(p^M,t^b)$ and after finite extension of coefficients.
\end{lemma}

\begin{proof}
The boundary has rank one over $L$, and the residual binary Shimura
datum is zero dimensional.  In its formal charts, a cusp tuple is
indexed by strictly positive integral $\alpha\in F$.  Here positivity
means positivity at both real embeddings.  There are no nonzero
rational indices on the boundary of the positive cone: an element
of $F$ with one embedding zero is zero.  Vanishing of the constant
term therefore imposes positive order on every relevant cone ray.

For $\alpha>0$, the coefficient is a section on the abelian
extension-data chart $C$ of
\[
                 \mathcal L_\alpha\otimes\mathcal E_\lambda|_C,
\]
where the Jacobi line bundle $\mathcal L_\alpha$ is ample.  This
description persists on the finite etale abelian covers in
\cref{lem:canonical-ordinary-level}.  The semiabelian differential
sequence has constant locally free graded pieces over the fixed
boundary base.  Apply \cref{lem:dual-fock-functor}, and tensor with
the determinant and opposite differential line factors.  This
gives $\mathcal E_\lambda|_C$ a finite filtration with constant
graded pieces, in precisely the integral lattice used by the theta
section.  This construction needs no rational splitting of the
semiabelian differential sequence.

We recall why the required vanishing is valid in every residue
characteristic.  If $L$ is ample on an abelian variety, choose
multiplication by an integer $n$ prime to the characteristic.  Its
trace splits $\mathcal O$ inside $[n]_*\mathcal O$, so
$H^r(L)$ is a summand of $H^r([n]^*L)$.  For sufficiently large such
$n$, the latter vanishes for $r>0$: $[n]^*L$ has arbitrarily large
ample degree, and Serre vanishing is uniform under the algebraically
trivial twists occurring in the theorem of the square, by properness
of the dual abelian variety.  The filtration just described gives
the same vanishing after tensoring with
$\mathcal E_\lambda|_C$.  Cohomology and base change, or induction
through the exact sequences for successive uniformizer powers,
then lifts every positive coefficient modulo $p^M$ integrally.

It remains to lift the arithmetic identifications between
coefficients.  Modulo translations, the stabilizer acts on positive
indices by unit norms.  Fixing a positive index forces this norm to
be one; the resulting CM norm-one unit group is finite.  The
remaining binary arithmetic stabilizer is finite by definiteness.
Neatness removes these finite stabilizers.  We may therefore choose
one lift in each free index orbit and transport it around the orbit,
without averaging or dividing by an orbit order.

The resulting tuples converge on every formal cone.  In fact, an
interior rational ray gives a linear functional strictly positive
on the nonnegative index cone.  A bound on its value bounds both
real coordinates of a positive index, and hence leaves only
finitely many points of the index lattice.  This is exactly the
formal convergence condition on the completed monoid algebra.
The formal-chart description and formal functions now identify
these tuples with the completed pushforward sections; see
\citet[Section~7.1.2]{Lan2013} and
\citet[Chapter~III, Theorem~11.1]{Hartshorne1977}.
Thus the required reduction map is surjective at a cusp.  Its
injectivity follows from applying the left-exact pushforward to
multiplication by a uniformizer on the torsion-free source sheaf.
Off the boundary the assertion is immediate from the finite level
normalization.  Faithfulness of completion for coherent modules
gives the assertion on the ordinary open.

The branch decomposition allows the zero lift on every other
branch.  Finally, $\cO[t]/(t^b)$ is a finite free $\cO$-module, so
the entire construction applies coefficientwise and gives the
stated Artin-ring assertion.
This proves the required special case directly; the general ordinary
subcanonical vanishing theorem is
\citet[Theorem~8.2.1.3, with the setup of Section~8.1]{Lan2018Ordinary}.
\end{proof}

\subsection{Hasse extension and the shifted weight}

The base-change lemma turns vanishing constant terms into a cusp
section modulo the working precision on the ordinary open.  We now
extend that section over the compactification and lift it to
characteristic zero, retaining its coefficients and Hecke relations
modulo the same precision.

Choose positive integers $m_i$ tending to infinity and tending
$p$-adically to zero on the fixed torsion branch, as in the preceding
sections.  Choose integers $M_i\to\infty$ sufficiently slowly along
the filter that, after the fixed clearing of denominators, all
constant coefficients vanish modulo $p^{M_i}$ through degree
$b-1$.  The theta section is then a cusp section on the ordinary
open over $\cO[t]/(p^{M_i},t^b)$ by
\cref{lem:cusp-base-change}.
We also require $M_i\leq n_i-c_b$, with $c_b$ a fixed constant
large enough for the truncated tame substitution of
\cref{sec:tame}.  Indeed, the coefficients of
$(1+t)^{p^{n_i}}-1$ below degree $b$ have valuations at least
$n_i-c_b$.  Thus the finite group relation is respected at this
precision.  All subsequent congruences use this one common
precision for the full ring.

The ordinary locus is the nonvanishing locus of the total Hasse
invariant.  A section of a coherent sheaf on this locus extends
over $X^{\min}$ after multiplication by a sufficiently high power
of Hasse.  The assertion follows from localization on a finite
affine cover.  It applies to the sheaf $\mathcal M_\lambda$ and to
each of the finitely many $t$-coefficients simultaneously.

For precision $M$, a sufficiently large $p$-power of Hasse lifts
modulo $p^M$.  To see this directly, choose local lifts of Hasse.
On overlaps their differences are divisible by $p$; raising to
$p^r$ makes those differences divisible by $p^{r+1}$, with the
bundle transition functions included.  They therefore glue for
$r$ sufficiently large.  In multiplicative differential frames,
the resulting lift reads $1$ modulo $p^M$.  Its exponent can be
increased arbitrarily while retaining this property.

Twisting farther by the ample Hodge line gives, by Serre vanishing
\citep[Chapter~III, Theorem~5.2 and Proposition~5.3]{Hartshorne1977}, surjectivity of
reduction on global sections.  We may consequently lift the
extended section to characteristic zero, coefficientwise in
$\cO[t]/(t^b)$.  Choose the Hasse weight shift $H_i$ arbitrarily
large and divisible by $(p-1)p^{M_i+c}$, with $c$ any needed fixed
constant.  The resulting weight, at each real place, is
\begin{equation}
 \lambda'_i
   =(m_i+H_i,\ m_i+H_i,\ 1+H_i;\ -1-H_i).
 \label{eq:shifted-weight}
\end{equation}
These generic sections are genuine cusp forms: the sheaf before
pushforward was twisted by minus the pulled-back boundary.

The size of $H_i$, the point at which Serre vanishing applies, and
the dimensions of the resulting spaces may depend on $i$, $m_i$
and $M_i$.  None of these choices consumes congruence precision.
If a finer neat cover was used, descent by its fixed averaging
operator costs only the valuation of its fixed index.  We absorb
this constant by reducing $M_i$ once.  There is no averaging at
the varying prime $q_i$.

\subsection{The two integral normalized operators}

Write $U_j$ for the actual double-coset operator, with the measures
and expanding-lattice convention used in the theta construction.
For either the original or shifted weight put
\begin{equation}
 \widetilde U_j=p^{-N_j(\lambda')}U_j,
 \qquad
 N_2(\lambda')=2+\lambda'_3+\lambda'_4,
 \qquad
 N_3(\lambda')=3+\lambda'_4.
 \label{eq:ordinary-normalization}
\end{equation}

\begin{lemma}[Integral trace calculation]
\label{lem:ordinary-trace}
The operators in \cref{eq:ordinary-normalization} preserve the
dual-Fock differential lattice on the canonical ordinary branch
and its reductions modulo every power of $p$.  After the Hasse
shift, both comparison eigenvalues are $1$ modulo the prescribed
precision.
\end{lemma}

\begin{proof}
It suffices to work at one distinguished $p$-adic factor.
The split PEL idempotents select the rank-three multiplicative
character lattice and the rank-one etale lattice; polarization
determines the conjugate factor.
Choose Serre--Tate coordinates $q_1,q_2,q_3$ for their bilinear pairing
\citep[Theorem~2.1(1),(2)]{Katz1981SerreTate}.  Locally
their deformation ring is
\[
 R=\cO[[q_1-1,q_2-1,q_3-1]].
\]
For $U_j$, expand the lattice transverse to a multiplicative
$j$-summand by $p^{-1}$.  The expanded multiplicative directions
are selected by $p^{j(3-j)}$ integral graphs.  Fix one graph.
Functoriality of the Serre--Tate pairing
\citep[Theorem~2.1(4)]{Katz1981SerreTate} for the lattice inclusion
and $p$ times the reverse inclusion shows that the remaining
extension choices replace the $j$ retained coordinates by their
$p$th roots; the other $3-j$ coordinates are unchanged.  Indeed,
expanding the etale generator requires a root when paired with a
retained multiplicative character, whereas the simultaneous
expansion of a transverse multiplicative direction cancels this
change for the other coordinates.

The corresponding finite free extension is the appropriate
completion of
\[
 R'=R[u_1,\ldots,u_j]/(u_a^p-q_a:1\leq a\leq j).
\]
Its basis consists of $\prod_a u_a^{e_a}$ with $0\leq e_a<p$.
The trace of every such basis element except $1$ is zero, and
$\Tr_{R'/R}(1)=p^j$.  Hence
\begin{equation}
                    \Tr_{R'/R}(R')\subseteq p^jR.
 \label{eq:root-trace-divisibility}
\end{equation}
The same monomial calculation applies on a boundary cone whenever
the Hecke correspondence takes a $p$th root of a torus coordinate.
For $j=2$ there are $p^2$ graphs and a degree-$p^2$ trace, and for
$j=3$ there is one graph and a degree-$p^3$ trace.  Thus the total
counts are $p^{j(4-j)}$, as required by the double cosets.

The transported multiplicative flags and differential comparison
maps depend only on the graph, not on these root choices.  They
can therefore be put outside the trace.  In an integral graph
basis, pullback scales $3-j$ positive differential directions by
$p$.  The opposite differential scales by $p^{-1}$, by duality
to the distinguished etale direction.

For the weights in use, the positive integral coefficient lattice is
\[
 (\det\omega^+)^{\lambda'_3}
 \otimes
 \mathcal D_{m-1}(\omega^+),
\]
and the opposite line occurs with exponent $-\lambda'_4$.
By \cref{lem:dual-fock-functor}, the integral graph changes preserve
this lattice.  For $j=2$, choose a graph basis in which the positive
differential map is $\diag(1,1,p)$.  It acts on the three minor
directions with entries $1,p,p$; hence on their monomial-dual
coefficient of degrees $(a_{12},a_{13},a_{23})$ it acts by
$p^{a_{13}+a_{23}}$.  This is integral for every $m$, without a
factorial or a change of lattice.  The determinant factor supplies
$p^{\lambda'_3}$.  For $j=3$, the positive differential map is the
identity in a graph basis.  The opposite line supplies
$p^{\lambda'_4}$ in both cases.  Before taking trace, the lower
bounds are consequently
\[
 \lambda'_3+\lambda'_4=N_2-2
       \quad(j=2),\qquad
 \lambda'_4=N_3-3
       \quad(j=3).
\]
Combining these with \cref{eq:root-trace-divisibility} gives
$p^{N_j}$ divisibility.  The graph sum introduces no denominator.
This proves integrality of $p^{-N_j}U_j$ on the stated lattice,
including when $N_j$ is negative.  The other $p$-adic factor is a
spectator, so the argument applies independently at both places.

In the original weight $(m,m,1;-1)$, both $N_j$ are $2$; the
theta eigenvalue $p^2$ thus gives normalized value $1$.
The shift $(H,H,H;-H)$ changes $N_2$ by zero and $N_3$ by $-H$.
Its differential determinant contributes $p^{(2-j)H}$ on the
same graph, precisely this change of normalization.  The remaining
unit frame factors, raised to the chosen divisible $H$, are $1$
modulo $p^M$; in canonical multiplicative frames the lifted Hasse
section itself is $1$.  Thus multiplication by the Hasse lift
intertwines the normalized operators modulo $p^M$.  Explicitly,
in the shifted weight
\[
                         N_2=2,\qquad N_3=2-H.
\]
Prime-to-$p$ Hecke correspondences preserve the canonical branch
and its integral lattice as well.  Their Hodge determinant changes
are $p$-adic units; the same divisibility of $H$ makes their
additional weight factors trivial modulo precision.  Hence all
prescribed comparison eigenvalues survive the lift.
\end{proof}

\subsection{The Hecke order and full Artin-ring factorization}

The lifted cusp section now has the theta comparison eigenrelations
on the ordinary locus.  To obtain a map from the Hecke order, we must
show that these values satisfy every relation in its classical action.
The nonconstant coefficient retained through the lift will let us
cancel the section from such a relation with only a fixed precision
loss.

Let $S_i$ be the classical cusp space in the weight given by
\cref{eq:shifted-weight} and at the prescribed finite level, with
the quotient detected on the canonical ordinary locus taken if
necessary.  We use its lattice of sections integral in the same
dual-Fock differential bundle, expressed in ordinary frames.
This is a full bounded lattice.  Indeed,
finitely many ordinary evaluations detect a finite-dimensional
classical space and bound the lattice above; clearing vertical
denominators of a basis bounds it below.  These statements may be
checked over a finite model field depending on $i$.
More explicitly, if $\mathcal Q_i^0$ is the module of integral
ordinary sections and $S_i$ is its rational classical subspace,
the lattice is $S_i\cap\mathcal Q_i^0$.  Consequently
$w\in S_i\cap p^M\mathcal Q_i^0$ implies
$w\in p^M(S_i\cap\mathcal Q_i^0)$.  Approximate classical
eigenrelations on the ordinary locus therefore give the same
precision in this lattice, without a further denominator loss.
\Cref{lem:ordinary-trace} and the prime-to-$p$ correspondences
preserve this lattice.

Define $T_i^0$ to be the $\cO$-algebra of commuting operators on
this space generated by the split good spherical coefficients
with Satake normalization $Q_v^{3/2}$, the operators
$\widetilde U_2,\widetilde U_3$, and the prescribed positive
binary controls with their Frobenius scalings.  Those latter
scalings use only powers of residue characteristics different
from $p$.  We retain the specified invertible center expansion
and last-block spherical degree-one control at each binary
place.  Over each finite model field the algebra is contained in
the endomorphism ring of the lattice.  It is therefore a finite
torsion-free $\cO$-algebra: an $\cO$-submodule of that finite
endomorphism lattice is finite.  Neither its rank nor the size
of the model field is required to be bounded as $i$ varies.

\begin{lemma}[Uniform cancellation in a truncated ring]
\label{lem:artin-content-cancellation}
Let $K'/k$ be a finite extension, and let
$a(t)\in\mathcal O_{K'}[t]/(t^b)$ satisfy
$v_p(a(0))\leq C$, where $C$ is a fixed nonnegative integer.
If $a(t)r(t)$ belongs to $p^M\mathcal O_{K'}[t]/(t^b)$, then
\[
       r(t)\in p^{M-bC}\mathcal O_{K'}[t]/(t^b).
\]
The bound is independent of the extension and of $a$.
\end{lemma}

\begin{proof}
Write $a=a_0(1+u)$, with $u\in t\mathcal O_{K'}[1/p][t]/(t^b)$.
Then
\[
 a^{-1}=a_0^{-1}\sum_{j=0}^{b-1}(-u)^j
          \in p^{-bC}\mathcal O_{K'}[t]/(t^b).
\]
Multiplication by this inverse proves the assertion.  The use of
the constant coefficient is essential, since it makes $a$
invertible over $K'[t]/(t^b)$.
\end{proof}

\begin{proposition}[Classical cusp congruence]
\label{prop:cusp-congruence}
After reducing $M_i$ by a fixed constant and relabeling it, there
are homomorphisms of $\cO$-algebras
\begin{equation}
 \lambda_i^0:T_i^0\longrightarrow
                  (\cO/p^{M_i}\cO)[t]/(t^{b_*}),
 \qquad M_i\longrightarrow\infty,
 \label{eq:hecke-congruence}
\end{equation}
which send every indicated operator to its theta comparison
eigenvalue, in the full truncated ring.  The weights are
\cref{eq:shifted-weight}, and the two normalized $p$-operators
have value $1$ at each distinguished place.
One may restrict $T_i^0$ to the local summand determined by
the residual augmentation system.  In that summand both
normalized $p$-operators are units.  No reduction of the operator
algebra is asserted here.
\end{proposition}

\begin{proof}
The comparison eigenvalues lie in the fixed coefficient field
$k$ and are integral in the specified normalizations.  This
follows from the theta parameter, the normalized value $1$ at
$p$, and the rational auxiliary block parameters.  Away from
$p$, the Frobenius-scaled polynomial coefficients are integral;
the character and residual-cardinality factors in these
normalizations are $p$-adic units.  The coefficients of
$(1+t)^\ell$ are integral for $\ell\in\Z_p$, and the preceding
choice $M_i\leq n_i-c_b$ makes the finite tame relations valid.
Consider the
polynomial $\cO$-algebra on all the chosen Hecke generators and
its evaluation at these comparison eigenvalues in the target of
\cref{eq:hecke-congruence}.  We must show that every relation
in its action on $S_i$ belongs to the kernel of this evaluation.

Let $s_i$ be the lifted cusp section, viewed with its $t$
coefficients.  On the canonical ordinary locus it agrees
modulo the working precision with Hasse times the integral
theta section.  Every generator acts there with the prescribed
comparison value modulo that precision, by
\cref{lem:ordinary-trace}.  The same is true for any integral
polynomial in the generators: multiplication and addition of
integral operators consume no precision, regardless of the
degree or number of its monomials.  If $P$ is a relation on the
classical cusp space, it follows that
\[
           P(\text{comparison values})\,s_i\equiv0
                                                   \pmod{p^{M_i}}.
\]
Evaluate the nonconstant coefficient furnished by
\cref{prop:bounded-content}, using the same Fock monomial and its
literal dual coefficient basis.  The definition of
$\mathcal E_{\lambda'_i}$ retains this integral pairing through the
lift.  Hasse is $1$ in that frame modulo precision, so its value
$a_i(t)$ has integral
coefficients and $v_p(a_i(0))\leq C$ for one fixed $C$.
The preceding display and
\cref{lem:artin-content-cancellation} imply
\[
 P(\text{comparison values})\equiv0
                                  \pmod{p^{M_i-bC},t^b}.
\]
The scalar polynomial on the left has coefficients in $k$;
integrality and divisibility can be descended from the finite
evaluation field because
$p^n\mathcal O_{K'}\cap\cO=p^n\cO$.
The fixed loss $bC$, together with the earlier fixed clearing
and neat-descent losses, is independent of $P$ and $i$.
Reducing and relabeling $M_i$ gives the factorization through
the entire operator order $T_i^0$.

Finally, a finite algebra over the complete local ring $\cO$
decomposes as a product of its factors at the residual maximal
ideals.  The target ring is local, with residue field obtained
by setting $t=0$ and reducing modulo the uniformizer.  Hence
$\lambda_i^0$ factors through the summand selected by this
residual augmentation.  The images of
$\widetilde U_2,\widetilde U_3$ are $1$, so these operators
lie outside that summand's maximal ideal and are units there.
This construction keeps every coefficient through degree
$b_*-1$ and permits nilpotents in $T_i^0$, as required for the
next section.
\end{proof}
